\documentclass[acmsmall,screen,review]{acmart}

\acmJournal{TOSEM}
\acmYear{2026}
\acmVolume{0}
\acmNumber{0}
\acmArticle{0}
\acmMonth{0}
\acmDOI{XXXXXXX.XXXXXXX}
\setcopyright{none}
\settopmatter{printacmref=true, printccs=false, printfolios=true}
\sloppy
\setlength{\emergencystretch}{3em}

\title{NOETHER: Data and Artifact Availability Statement}

\newcommand{\threeaffil}{%
  \affiliation{%
    \institution{School of Computing, University of South China}
    \city{Hengyang}
    \postcode{421001}
    \country{China}
  }%
  \affiliation{%
    \institution{Hunan Engineering Research Center of Software Evaluation and Testing for Intellectual Equipment}
    \city{Hengyang}
    \postcode{421001}
    \country{China}
  }%
  \affiliation{%
    \institution{CNNC Key Laboratory on High Trusted Computing}
    \city{Hengyang}
    \postcode{421001}
    \country{China}
  }%
}

\author{Meng Li}
\authornote{Corresponding author.}
\orcid{0000-0002-1074-1502}
\email{mlemon@usc.edu.cn}
\threeaffil

\author{Xiaohua Yang}
\orcid{0000-0002-2977-1787}
\email{xiaohua1963@foxmail.com}
\threeaffil

\author{Jie Liu}
\orcid{0009-0008-1970-8347}
\email{jieliu@usc.edu.cn}
\threeaffil

\author{Shiyu Yan}
\orcid{0000-0001-7626-5185}
\email{yanshiyu@usc.edu.cn}
\threeaffil

\begin{document}
\maketitle

\section*{Data and Artifact Availability}

To support independent verification, the supplementary code and data are archived
on Zenodo under DOI
\url{https://doi.org/10.5281/zenodo.20250634}.
The same materials are supplied through the submission system with the same file
manifest and content hash.

\paragraph{Review-stage supplementary archive.} The following supplementary
materials are submitted alongside the manuscript and are available through the
submission system under SHA-256 content hash.
The hash is computed over the concatenated tar-archive of the listed items and
reported in the final manuscript at acceptance.
\begin{itemize}
    \item \textbf{S1} Python reference implementation of CONSTRUCT-MP described
    in supplementary~S9, including the Boltzmann and equivariant-ML
    instantiations.
    \item \textbf{S2} The full PWR-MR corpus, with each MR annotated by canonical
    structural component, NOETHER MetaPattern, and source equation.
    \item \textbf{S3} The SE(3)-equivariant point-cloud testing harness,
    including the mutation generators and MR sets used in the case-study
    comparison.
    \item \textbf{S4} A reproducibility manifest documenting the random seeds,
    model checkpoints, and dataset versions referenced in the case study.
    \item \textbf{S5} GenMorph/PIT pilot material, head-to-head summaries,
    parsed mutation results, and statistical scripts for the secondary
    executability and complementarity analyses.
    \item \textbf{S6} Query-optimiser pilot material for the
    relational-equivalence extension, including sampled Calcite/QED plan pairs
    and oracle scripts.
    \item \textbf{S7} Defects4J algebra-rich substrate material,
    pre-registration records, SUT-selection criteria, and
    $\mathcal{L}^{*}$-blindness evidence.
    \item \textbf{S8} METRIC+ comparison material, Sun et al. subject
    re-implementations, scope analysis, and head-to-head protocol/results files.
    \item \textbf{S9} Appendices and detailed empirical tables migrated out of
    the Fast Impact main body for length conformance.
    \item \textbf{S10} Home-field generation-to-detection checks,
    rank/faithfulness scripts, and additional Invariance-Blindness sanity
    checks.
    \item \textbf{S11} Industrial N5 coverage classification inputs and results.
    \item \textbf{S12} Cross-domain N5 numerical-library and signal-processing
    derivation checks and outputs.
\end{itemize}

\paragraph{Public archival release.} The public archival package is deposited on
Zenodo under DOI
\url{https://doi.org/10.5281/zenodo.20250634}. The
camera-ready version will additionally record the final SHA-256 hash of the
submission archive. We aim for the \emph{Available} and \emph{Functional}
artefact-evaluation badges.

\end{document}
